\section{Design}
\label{sec:design}

\sys is two mechanisms and no controller. This section takes the five problems
of \S\ref{sec:problems} in order (Table~\ref{tab:map}) and says what addresses
each, including the one we tried first and discarded.

\begin{table}[t]
\centering
\small
\begin{tabular}{@{}llp{3.1cm}@{}}
\toprule
 & problem & addressed by \\
\midrule
P1 & burst of unknown size & \S\ref{sec:d-burst}: discard, not growth \\
P2 & the sum is global & \S\ref{sec:d-cap}: cache-derived cap \\
P3 & larger is worse past the ceiling & \S\ref{sec:d-burst}, and why we do not size \\
P4 & work spent before the drop & \S\ref{sec:d-place}: discard in the driver \\
P5 & per socket vs per queue & \S\ref{sec:d-select}: select by L4 type \\
\bottomrule
\end{tabular}
\caption{Each problem and the part of the design that answers it.}
\label{tab:map}
\end{table}

\subsection{P1 and P3: absorb the burst by not doing the work}
\label{sec:d-burst}

P1 asks for a buffer large enough to hold one moderation window. P3 says a
larger buffer is harmful past the point where the consumer falls behind. Taken
together they rule out the obvious answer, because the two are not separate
regimes: the same socket crosses from one to the other as its offered rate
rises, and a mechanism that grows the buffer for P1 is still growing it when P3
applies.

\sys therefore does not grow the buffer. When the socket queue is full, the
driver stops delivering UDP completions for a short window, discarding them
before an \texttt{skb} exists. The burst is absorbed by shrinking it rather than
by making room for it, and the buffer stays small enough that what is queued is
still resident when the consumer arrives.

\paragraph{The window is fixed at 50\,\textmu s.} We swept it from 25 to
1600\,\textmu s and found the response asymmetric: too short costs 7\%, too long
costs 67\%, and 1600\,\textmu s is worse under a load transition than no window
at all. An adaptive controller is possible; the oracle's advantage over a fixed
50\,\textmu s is 1.6\%, which is less than the cost of searching for it.

\paragraph{Why not size the buffer.} We built the adaptive sizer the problem
statement suggests---TCP's rule with the interval between the receive queue
filling and emptying standing in for the RTT---and report it in
\S\ref{sec:sizing}. It works, and it is not better than this. The reason is P3:
the estimate it converges on is ``as much buffer as this socket can use'', and
the measurement says the quantity worth minimising is how much data sits between
producer and consumer.

\subsection{P2: a cap derived from the cache}
\label{sec:d-cap}

\sys reads the last-level cache size from the kernel's cache topology at boot
and takes a configurable fraction of it, one half by default, as an allowance. A
socket's whole buffer is charged against the total once, on the first packet it
receives, and returned when the socket is destroyed.

\paragraph{A bound, not a target.} The allowance does not look for a good size.
It prevents individually reasonable sizes from summing to a bad one, which is
exactly what P2 says a per-socket limit cannot do. Within the allowance, what
the application or the operator asked for stands.

This is a weaker mechanism than a controller, and deliberately so.
\S\ref{sec:workingset} establishes a claim about the sum and says nothing about
where within the bound a given socket should sit. A cap is the mechanism that
claim licenses.

\paragraph{What the implemented cap does and does not do.} It refuses any
\emph{increase} that would take the total past the allowance. It does not
currently refuse the buffer a socket starts with: a machine configured with an
8\,MB \texttt{rmem\_default} and eight sockets is charged 64\,MB and permitted
it. The measurement in Table~\ref{tab:nostatic} showing that discarding fails to
rescue eight 8\,MB buffers is therefore evidence that such a cap is
\emph{needed}, not that ours enforces it. Clamping the default---as against an
explicit \texttt{SO\_RCVBUF}, which we argue below should be honoured---is the
obvious extension; it requires deciding what a socket gets when the allowance is
already spent, and we have not measured our way to an answer.

\paragraph{Capacity, not occupancy.} The allowance bounds granted capacity
rather than queued bytes. Bounding queued bytes is the natural first attempt and
does not work: while the consumer keeps up the queues sit nearly empty---0.36\,MB
held against 8\,MB of capacity in one measurement---so the bound does not bite
until the queues have grown deep, by which point the buffers behind them are far
too large. Occupancy follows capacity, not the reverse.

\paragraph{All of it is charged.} The whole buffer, not the part above some
default. Counting only the excess leaves the baseline invisible, and the
baseline is not small when there are many sockets: eight at a 1\,MB default are
8\,MB of a 9\,MB allowance before anything else happens. An earlier version
charged only what it had itself handed out, and eight sockets reached 17.8\,MB
against a 9\,MB allowance.

\paragraph{Sockets that set \texttt{SO\_RCVBUF}.} An application that names a
size still receives it; overriding an explicit request would break a long
standing expectation, and in isolation the request is rarely the problem. \sys
charges such a socket for what it holds so that everyone else stops compounding
it: 64\,MB is entirely safe until it is the eighth one.

\paragraph{Why the fraction is an operator setting.} We would prefer the
allowance to be fully derived, and it nearly is---the cache size comes from the
kernel's own topology, so it follows the hardware. What the kernel cannot see
from the socket layer is how much of that cache is already committed. The
dominant unseen term is the NIC's descriptor pages, 16\,MB for a 1024-entry ring
on our hardware, which the driver allocates and never surfaces in a form the
socket layer can query. A fully automatic version is possible---the driver knows
the ring size and the descriptor format---but it would put the mechanism behind
a per-driver change, and one number per machine is a considerably better
interface than one number per socket.

\subsection{P4: discard where the work has not yet been done}
\label{sec:d-place}

Placement is the whole of P4. A drop at the socket happens after the packet has
cost descriptor processing, an \texttt{skb} allocation, GRO, and a route lookup;
a drop at the completion queue costs a comparison. \sys marks the receive queue
on which a socket drop occurred, and for the window the driver discards matching
completions in \texttt{mlx5e\_handle\_rx\_cqe()} before any of that work
happens.

\paragraph{Take the queue from the packet, not the socket.} The queue to hold
back is the one that delivered the packet which could not be queued. Reading it
from the socket instead, via \texttt{sk\_rx\_queue\_get()}, is wrong in a way
that appears only with more than one queue: that field is maintained for
connected sockets, and an unconnected socket---the many-senders case UDP exists
for---reads back $-1$, which a fallback turns into queue zero. Every such socket
then throttles queue zero regardless of which queue was drowning.

\subsection{P5: discard selectively}
\label{sec:d-select}

The completion descriptor carries the L4 header type, so the discard can be
restricted to UDP. This is what allocation structurally cannot do. Buffers are
per socket while the queue is shared, so when TCP and UDP contend for one queue,
no setting of the UDP socket's buffer gives TCP more of it; what TCP is short of
is service, not space. Declining to build \texttt{skb}s for one protocol leaves
the core free for the other, and \S\ref{sec:evalproto} measures the result.
